\cotefile{15}
\scanpage{195}
2°) $X_{\deleted{\unsure{\text{réd}|\text{red}}}}$ domine $\operatorname{Spec}\mathbb{Z}$, i.e. $X$ contient des pts de car résiduelle nulle.

Alors \deleted{en plus de ces \inserted{dispose \deleted{\unsure{des}}} (et essentiellement \unsure{(i)s|(i)} (iii) dans le cas d’inégales car) on a, pour tt pt} \deleted{\unsure{(iii)} (iii) bis)}

(i) $\unsure{\xi|\zeta}_{\ell}\in \unsure{\mathbb{T}|T|\mathbf{T}}_{\unsure{\alpha|x}}(\mathbb{Z}_{\ell})$ \uline{tt} $\ell\in\mathbb{P}$, $\xi_\ell$ provenant sans doute de $T_\alpha(\widetilde{\mathbb{Z}}_{\ell\mathbb{Z}})$

$\xi_\ell$ dépend du choix d’un pt géom quelconque de $X$, et est indép. de ce choix modulo opération de $\unsure{\Pi|\pi}_1$, notation $\xi_{\bar{x},\ell}$ ou $\xi_{x,\ell}$

(ii) Si $X$ a des points $x$ de car. $p>0$, alors on dispose également des $\xi_{x,p}$, $\xi_{x,p}^{0}$.

\unsure{Correction|Connexion|Connection|Conjecture} : relation, \deleted{\unsure{de la}} si $y$ est un pt de car. $p>0$ et \deleted{\unsure{x}} un pt de car. $0$, entre $\xi_{y,p}$ et $\xi_{x,p}\otimes_{\mathbb{Z}_p}\mathrm{W}(k)$ [$y$ pt à val. ds corps parfait $k$]

\uline{Ils ne sont en général pas isomorphes}, sans doute. \deleted{\unsure{Ils}} Si \deleted{\unsure{ne sont pas}} $k$ est fini, ils sont cependant \inserted{isomorphes \unsure{par Lang|par Lang,}} \deleted{\unsure{et les groupes devant}} \deleted{\unsure{groupes de ces objets}} \deleted{\unsure{devant}} \deleted{\illegible{1}}

(iii) Si $x$ est un pt à val. dans un corps \uline{de car $0$}, on a $\xi_{x,\infty}\in T_\alpha(k)$ par Hodge-de Rham. \inserted{Avec ses \textcircled{\uuline{2}} filtrations \deleted{\unsure{et le groupe}}} Il n’y a pas lieu de

% margin note written sideways
\margin{L’une est celle qui donne comme gradué associé les $\unsure{H^{q}|H^{i}|H^{\bullet}}(x,\Omega^{\unsure{p|j|P}})$, l’autre, (\unsure{arithmétique|authentique|carthentique}) fait pendant à la filtration canonique des motifs. (Elle est \uline{triviale} pour les motifs simples.) La première filtration est originale au foncteur Hodge-De Rham, bien faire attention qu’elle ne provient pas d’une filtration du motif, c’est une filtration du \uline{foncteur}\ldots. !}

% archive sheet number: 194
\scanpage{196}
ALGÈBRES de LIE SEMI-SIMPLES \mspage{(7)}

Tribu n° 59

mettre l'exer. 1, (p. 36), en c) de la prop. 11, et l'exer. 2 dans le texte.

Ajouter un utile rablot sur les algèbres de la forme, (commutative), $\mathrm{x}\ \underline{\mathrm{sl}}_2$, à savoir les généralisations des prop. 5 et 6.

\S3 - \uuline{Le système de racines}\ldots

\uline{n° 1} - Définir ``déploiement'', (i.e. couple $(\underline{\mathrm{g}},\underline{\mathrm{h}})$), ``déployable''. Dire que $(*)$ de la déf. 1 signifie que $\operatorname{ad}(\mathrm{x})$ est diagonalisable, (on sait qu'il est semi-simple). Ajouter ``semi-simple'' à la l.2 bas, (p. 37).

Remplacer ``canonique'' par ``diagonale'', et ``abélienne'' par ``commutative'', ousqui faut.

\uline{n° 2} - Il s'agit du sous-espace ``primaire'' à la p. 38, l.4. A la l.5 dire : ``comme $\operatorname{ad}(\underline{\mathrm{h}})$ est semi-simple, $\underline{\mathrm{g}}^\alpha$ est l'ensemble\ldots'', et ne pas répéter ça à la l.16. Dans l'énoncé du th. 2 dire que $\underline{\mathrm{h}}^\alpha$ est un sous-espace de dim. 1 de $\underline{\mathrm{h}}$. Y expliciter l'isom. $\varphi_\alpha$ dans (iii). Ajouter ``et si'' dans le (iv).

\uline{n° 3} - Après le premier alinéa de la p. 41, démontrer le th. 3, c), (cf. p. 44), ce qui règle son compte au th. 3. Plus bas utiliser l'isom. $\varphi_\alpha$. Commencer la p. 42 par ``Soit\ldots''. A la l.3, (p. 42), qualifier $\mathrm{W}(2)$ de ``module simple''. Ce qui est entre parenthèses en-dessous est maintenant fait au \S\ 2. Dire, (p. 42, l.12), que $\mathrm{R}$ est un système de racines `` de $\underline{\mathrm{h}}$ ''.

Le haut de la p. 44 remonte p. 41.

\uline{n° 4} - Titre : terminologie et inventaire.

\uline{n° 5} - A la l.2 bas de la p. 45 dire : ``Dans tout couple $(\alpha,-\alpha)$ de racines opposées, choisissons arbitrairement une racine $\alpha$ ; pour chaque racine $\alpha$ ainsi choisie, choisissons\ldots\ $\mathrm{X}_\alpha$\ldots''.

% archive sheet number: 195
\scanpage{197}
penser que pour des pts différents à val. dans le \unsure{$=$|même} $k$, les deux foncteurs obtenus, soient isom (pas plus que dans le cas (i)). Il n’y a plus en général d’opérations de $\unsure{\Pi|\pi}_1(X)$ sur $\unsure{\xi|\zeta}_{x,\infty}$, mais encore comme dans (i) les autom de $k$ sur $k(x)$ opèrent dans $\xi_{x,\infty}$ $\tfrac12$ \uline{\unsure{linéaimt|linéairemt|linéairement}} (donnée de descente). Un autom$^{t}$ de $k$ qui n’est pas l’identité sur $k$ définit $\xi_{x,\infty}^{\sigma}$ qui à priori n’est plus isomorphe à $\xi_{x,\infty}$.

Un cas intéressant est celui où $k=\mathbb{Q}$, on trouve alors une \deleted{\unsure{trivialit.|trivial}} $\xi_{x\infty}\in T_{\unsure{\alpha|x}}(\mathbb{Q})$, i.e une trivialisation de $T_{\mathbb{Q}}$. Deux trivialisations à val dans $T_{\mathbb{Q}}$ ne sont en général pas isom à priori \unsure{les mêmes}. De plus, il n’y a pas d’isom \uline{canonique} prévisible entre

\[
\xi_{x,\infty}\ \text{et}\ \xi_{x\infty}\otimes_{\mathbb{Q}}\mathbb{Q}_{\ell}\ \text{et}\ \xi_{\bar{x},\ell}\otimes_{\mathbb{Z}_{\ell}}\mathbb{Q}_{\ell}\quad -
\]

\deleted{bien que pas} et comme le premier à priori ne provient pas d’un objet de $T_{\unsure{-|\alpha}}(\mathbb{Z}_\ell)$ mais seult de $T(\mathbb{Q}_\ell)$, il n’y a pas lieu de penser qu’il y ait $=$ un isom. Ainsi, tt pt rat / $\mathbb{Q}$ définit une classe de coh dans les $H^{2}(\mathbb{Q}_\ell,\unsure{\mathbb{G}|G}^{\circ(\alpha)}_{\ell})$. Dans le \deleted{\unsure{de}} cas général, il faut remplacer $\mathbb{Q}_\ell$

% slanted
\margin{en effet, ils sont déjà non isom pour les Motifs de poids et niveau $0$}

% archive sheet number: 196
\scanpage{198}
ALGÈBRES de LIE SEMI-SIMPLES \mspage{(6)}

Tribu n° 59

Un cor. à la prop. 1 dira : ou bien tous les $\mathrm{e}_{\mathrm n}$ sont lint. indép. (car de poids différents), ou bien l'un est nul et alors $\lambda$ est un entier $\mathrm{q}\geqslant 0$, indépendance jusqu'à $\mathrm{e}_{\mathrm q}$, nuls ensuite. Ceci remplace une partie de la prop. 2. Dans l'exer. du bas de la p. 26, ajouter la stabilité par les binomiaux $\binom{\mathrm{H}}{\mathrm{n}}$.

\uline{n° 2} - Titre : les modules simples $\mathrm{W}(\mathrm{m})$, (dire ``simple'' au lieu de ``irréductible''). Dans la prop. 3 convenir que $\mathrm{e}_{-1}=0$ et $\mathrm{e}_{\mathrm n}=0$ pour $\mathrm{n}\geqslant\mathrm{m}+1$. Prendre $\mathrm{e}=\alpha^{\mathrm m}/\mathrm{m}!$ ; alors $\mathrm{e}_{\mathrm n}=\mathrm{Y}^{\mathrm n}\mathrm{e}/\mathrm{n}!$ ; qualifier la base $(\mathrm{e}_{\mathrm n})$ de ``canonique'' ; arranger (4).

Remplacer l'exer. 2, (p. 28), par un truc de Gabriel. Mettre l'exer. 3, (p. 29), dans le texte. Remonter ici le lemme 4 des pp. 53-54.

\uline{n° 3} - Titre : décomposition des $\underline{\mathrm{s}}$-modules. Ajouter ``de dim. finie'' dans la prop. 5 et son cor. Il semble qu'on n'a pas besoin du (ii) et du (iii) de la prop. 6 ; si oui, ça passe en exer. et on bloque le (i) avec la prop. 5. Dire ``dim. finie'' dans la prop. 7.

A la p. 30, l.12, donner une référence pour la complète réductibilité. Mieux détailler la remarque 1), p. 31. A la place de la remarque 2), (pas si ``facile''), on donnera le calcul du Casimir, (cf. Papa Cartan).

\uline{n° 4} - Dire, (l.5, p. 31), qu'on a $\mathrm{x}^2=0$. Dire que $\mathrm{E}$ est de dim. finie dans la prop. 8.On préfère la seconde démonstration p. 34, (les exponentielles de loct. nilpotents seront définies en Alg. IV). A la fin de la remarque, (p. 34 bas), dire : ceci montre aussi que ce sont les seules représentations irréd. holomorphes de $\mathrm{G}$. Expliciter la fin de la démonstration de la prop. 9. Lire `` $\mathrm{E}=\mathrm{S}^{\mathrm m}(\mathrm{V})$ '' à la l.4 bas de la p. 35. Dans la prop. 11, a) dire que $\unsure{\vartheta|\theta}^4=1$ ;

% archive sheet number: 197
\scanpage{199}
par $\mathbb{Q}_{\ell}(\unsure{k|K})$ --- c’est particulièrement intéressant dans le cas $k$ un corps de nbs.

(iv) Si $x$ est un pt à valeurs \uline{dans $\mathbb{C}$}, il lui correspond

\[
\unsure{\xi|\zeta}_{x,\mathbb{Z}}\in\mathbb{T}_{\unsure{\alpha|x}}(\mathbb{Z})\qquad
\xi_{\alpha}(\mathbb{Z})(F)=\text{réseau entier des }F_x,
\]

\deleted{\unsure{fonct|Quand} --- De plus, le groupe $\mathbb{Z}/2\mathbb{Z}$ opère sur $\mathbb{C}$ \inserted{$\xi_{x,\mathbb{Z}}$ et $\xi_{\unsure{\tau x|x},\mathbb{Z}}$ sont \unsure{can} isom.} \unsure{foncteurs par $\mathbb{Z}$} ($\tau$ est la conj. complexe),}

et des isom can

\[
\left\{\begin{aligned}
\xi_{x,\mathbb{Z}}\otimes_{\mathbb{Z}}\mathbb{Z}_{\ell}&\simeq\xi_{x,\ell}\\
\xi_{x,\mathbb{Z}}\otimes_{\mathbb{Z}}\mathbb{C}&\simeq\xi_{x,\infty}
\end{aligned}\right.
\]

D’ailleurs $\xi_{x,\mathbb{Z}}$ est défini en termes d’un pt à val. dans un corps \inserted{top} $K$ isom à $\mathbb{C}$ (à isom \unsure{can} près) en d’autres termes on a un isom canonique ``involutif''

\[
\xi_{x,\mathbb{Z}}\simeq\xi_{\tau x,\mathbb{Z}}
\]

où $\tau x$ est le pt complexe conjugué.

\deleted{Donc} On peut dire aussi, si $x$ est localisé en $a\in X$, que \inserted{a) si $k(a)$ est dense dans $\mathbb{C}$, alors} $\xi_{x,\mathbb{Z}}$ dépend de la donnée d’une topologie \inserted{archimédienne} (ou valuation) sur $k(a)$, à complété $\mathbb{C}$ \enspace b) si $k(a)$ \unsure{non dense dans|non dense ds}.

% archive sheet number: 198
\scanpage{200}
ALGÈBRES de LIE SEMI-SIMPLES \mspage{(5)}

Tribu n° \unsure{59|57|53}

\uline{n° actuellement 4} - Dire ``produit fini'' dans la déf. 3. Dire que $\operatorname{Autel}(\mathrm{g})$ est invariant dans $\operatorname{Aut}(\mathrm{g})$. Au milieu de la p. 12 raisonner sur $\underline{\mathrm{g}}/\underline{\mathrm{h}}$, (au lieu de Fitting).

Ajouter `` $(\mathrm{t}\in\mathrm{k})$ '' à la l.12 p. 13 ; dire ``une seule valeur'' à la l.14. Il s'agit d'une inclusion à la l.18. Dire `` image par $\mathrm{F}$ '' à la l. 19.

Dans l'énoncé du th. 2 dire ``le groupe $\operatorname{Autel}(\mathrm{G})$'' opère transitivement'' sur les Cartans. Dire ``la conclusion de (iii)'' dans la remarque p. 14.

Ajouter le cor. suivant : soient $\underline{\mathrm{h}}\subset\underline{\mathrm{g}}'\subset\underline{\mathrm{g}}$ ; si $\underline{\mathrm{g}}'$ a même rang que $\underline{\mathrm{g}}$ et contient un élément régulier de $\underline{\mathrm{g}}$, alors $\underline{\mathrm{h}}$ est une Cartan de $\underline{\mathrm{g}}$ si et seult. si c'est une Cartan de $\underline{\mathrm{g}}'$, (utiliser la prop. ajoutée au n° 3).

\uline{n° 6} - Dire partout ``commutative'', (au lieu de ``abélienne''). Ajouter ``car nilpotente'' à la l.9, p. 16. Utiliser la prop. ajoutée au n° 3 dans la prop. 14.

\uline{n° 7} - Style du th. 4 : ``le groupe des\ldots\ opère transitivement\ldots''. Idem dans le cor. p. 20. Dire ``composante neutre'' à la p. 21, l.3. Ajouter que, si $\underline{\mathrm{g}}$ est une alg. Lie sur $\mathrm{K}$, et si $\mathrm{k}$ est un sous-corps de $\mathrm{K}$, les Cartans de $\underline{\mathrm{g}}$ considérée comme algèbre sur $\mathrm{k}$ sont les mêmes que celles de $\underline{\mathrm{g}}$ considérée comme algèbre sur $\mathrm{K}$.

\S2 - \uuline{Les représentations de} $\mathrm{sl}_2$.

Eviter les points, (p. 24). Dire : ``sous-algèbre diagonale''.

\uline{n° 1} - Dans la remarque, (p. 25), dire ``(car $\mathrm{X}$ est un crochet)''. Généraliser la déf. des poids à tous les vecteurs propres de $\mathrm{H}$. Un exer. montrera les séries génératrices correspondant à (2). A la p. 25, l.2 bas, dire `` d'où, en divisant par $\mathrm{n}!$ ''.

% archive sheet number: 199
\scanpage{201}
$\mathbb{C}$, i.e. $\overline{k(a)}\simeq\mathbb{R}$\inserted{isomor. can.}, alors $\unsure{\xi|\zeta}_{x,\mathbb{Z}}$ est encore défini par la valuation archimédienne sur $k(a)$, mais \deleted{suivant que les} \inserted{\unsure{conformt aux|conformément aux|compte tenu}} choix possibles dans la clôture algébrique de $k(a)$, le gpe de Galois de $\mathbb{C}=\overline{\mathbb{R}}$ sur $\mathbb{R}$ opère dans $\xi_{x,\mathbb{Z}}$, soit \uline{$\tau$ cette opération}, qui est involutive. [On ne la confondra pas avec l’\unsure{endom|endom.|action} \uline{$\tfrac12$ linéaire} de $\xi_{\bar{a},\infty}$ correspondant à $\tau$ ; il n’y est pas question par exemple d’un réseau stable].

\margin{On écrit aussi $\xi_{a,\mathbb{Z}}$ au lieu de $\xi_{x,\mathbb{Z}}$,}

% slanted
\margin{on choisit le \uline{\unsure{plon|plus|plgt}} défini par le pt à val dans $\mathbb{C}$}

Dans ce cas, les invariants sont

\diagram{\begin{tikzcd}[ampersand replacement=\&,column sep=small,row sep=small]
\xi_{a,\infty}\in T_{\alpha}(\mathbb{R}) \arrow[dr,dash] \& {}\\
\deleted{\xi_{\unsure{a|\bar a},\infty}\otimes_{\mathbb{R}}\mathbb{C}\simeq} \& \xi_{\bar{a},\infty}\in T_{\alpha}(\mathbb{C})\\
\xi_{\bar{a},\mathbb{Z}}\in T_{\alpha}(\mathbb{Z}) \arrow[ur,dash] \& {}
\end{tikzcd}}

\[
\xi_{\bar{a},\infty}\overset{\text{\uuline{can}}}{\simeq}\xi_{a,\infty}\otimes_{\mathbb{R}}\mathbb{C}\simeq\xi_{\bar{a},\mathbb{Z}}\otimes_{\mathbb{Z}}\mathbb{C}
\]

$\unsure{\Pi|\pi}_1(X)$ n’opère pas à priori dans $\xi_{a,\mathbb{Z}}$, les seules opérations sont celles de $\tau$ si $a$ est \uline{réel}. Cependant, sur $\xi_{\bar{a},\infty}$, il y a (en plus éventuellement de $\tau$, si $a$ est réel) les opérations de $\operatorname{Aut}_{k(a)}^{\unsure{\ell|i}}(k)$, qui sont $k$ linéaires.

% archive sheet number: 200
\scanpage{202}
ALGÈBRES de LIE SEMI-SIMPLES \mspage{(4)}

Tribu n° 59

multiplicité $\geqslant\mathrm{n}$. On donnera le principe de Lefschetz comme lemme. En cor. : si $\mathrm{G}$ est connexe, (sur $\mathrm{R}$ ou $\mathrm{C}$), on a $\mathrm{U}=\mathrm{G}$.

\uline{n° 2} - Remplacer ``canonique'' par ``diagonale'' à la l.5, p. 7. Dans la dém. de la prop. 6 remplacer ``la remarque faite ci-dessus'' par ``l'exemple 3)'', et ajouter à la fin ``d'après l'exemple 1)''. Dire que toute sous-algèbre qui contient une Cartan est son propre normalisateur.

Dans la prop. 8, c'est au chap. I, p. 53 qu'il faut référer.

\uline{n° 3} - Dire, au début, qu'on pose $\mathrm{a}_{\mathrm{n}}(\mathrm{x})=1$. Dire que le rang est la multiplicité de la racine $0$. Dans l'ex. 1, ajouter que, si $\underline{\mathrm{g}}$ est nilpotente, le rang est la dimension, (cf. prop. 1). Vérifier la formule du bas de la p. 9. Les deux lignes du haut de la p. 10 ne sont vraies qu'en caract. $\neq 2$, (en caract. 2, $\underline{\mathrm{sl}}_2$ est nilpotent). L'ex. 3), (haut p. 10), est trop sec.

Dans le th. 1, il est plus simple de raisonner sur $\underline{\mathrm{g}}/\underline{\mathrm{h}}$, (au lieu de Fitting). Dire ``sous-algèbre de Cartan'' dans le cor. Vider la remarque finale, (fausse en caract. $\mathrm{p}$).

Ajouter la prop. suivante : soient $\underline{\mathrm{h}}\subset\underline{\mathrm{g}}'\subset\underline{\mathrm{g}}$ des alg. Lie ; alors a) si $\underline{\mathrm{h}}$ est une Cartan de $\underline{\mathrm{g}}$, c'est une Cartan de $\underline{\mathrm{g}}'$, et tout élément de $\underline{\mathrm{g}}'$ régulier dans $\underline{\mathrm{g}}$ est régulier dans $\underline{\mathrm{g}}'$ ; b) si $\underline{\mathrm{g}}'$ est le nilespace d'un nilpotent, alors tout élément régulier de $\underline{\mathrm{g}}'$ est régulier dans $\underline{\mathrm{g}}$, et a mêmes nilespaces dans $\underline{\mathrm{g}}$ et $\underline{\mathrm{g}}'$ ; de plus $\underline{\mathrm{h}}$ est une Cartan de $\underline{\mathrm{g}}$ et $\underline{\mathrm{g}}'$.

\uline{n° actuellement 5} - L'assertion de la l.6, p. 15, est fausse en caract. $\mathrm{p}$ ; remplacer par un petit calcul bébête sur la multiplicité de la valeur propre $0$, (cf. remarque 2).

% archive sheet number: 201
\scanpage{203}
En plus de la \uline{structure entière} sur le \deleted{\unsure{él}} \inserted{foncteur} de Hodge-DeRham, relatif à un pt complexe, il faut signaler la \uline{décomposition en types}, également de nature transcendante (elle est compatible avec la filtration canonique). On peut aussi dire que cette dernière \deleted{\unsure{associée}} \inserted{est déterminée via filtration} par \uline{l’opération involutive} \inserted{$\tau$ $\tfrac12$ linéaire} $\tau'$ sur $\unsure{\xi|\zeta}_{x,\infty}$, provenant de l’écriture

\[
\xi_{x,\infty}=\xi_{x,\mathbb{R}}\otimes_{\mathbb{R}}\mathbb{C}
\qquad\text{où}\quad \xi_{x,\mathbb{R}}\simeq\xi_{x,\mathbb{Z}}\otimes_{\mathbb{Z}}\mathbb{R}.
\]

Si $k(a)$ est réel, on a une autre application $\tfrac12$ linéaire involutive $\tau''$, provenant de

\[
\xi_{\unsure{\alpha|x},\infty}\simeq\xi_{a,\infty}{}_{\mathbb{R}}\otimes_{\mathbb{R}}\mathbb{C}
\]

et $\tau$ sur $(\xi_{x,\mathbb{Z}})_{\mathbb{C}}$ n’est autre que $\tau'\tau''$

\diagram{\begin{tikzcd}[ampersand replacement=\&,column sep=0.3em,row sep=1.2ex]
{} \& \tau={} \& \tau' \arrow[dl,dash,bend right=25] \& \tau'' \arrow[dr,dash,bend right=20] \& {}=\tau''\tau' \\
{} \& \substack{\tfrac12\ \text{linéaire}\\\text{transcendant}} \& {} \& {} \& \substack{\tfrac12\ \text{linéaire}\\\text{algébrique}}
\end{tikzcd}}

% slanted
\margin{à vérifier}

La situation locale \inserted{(dans le cas réel)} est connue par la connaissance de $\xi_{\unsure{x|a},\infty}^{\tau''}=(\xi_{a,\infty})_{\mathbb{R}}$ \inserted{(Hodge-DeRham réel)}, muni de $\tau=\tau'$ (involution transcendante, qui est $\mathbb{R}$-linéaire sur le dit espace), ou \inserted{mieux par celle de $\tau=\tau''$ opérant sur $\xi_{x,\infty}^{\tau'}\simeq\xi_{x,\mathbb{R}}$,} \uline{\unsure{évidemt} de $\tau$ sur $\xi_{a,\mathbb{Z}}$}, \unsure{donc au choix de|d’où on déduit de|dont on déduit de}

\[
\boxed{\tau\ \text{sur}\ \xi_{a,\ell}}
\]

[or c’est là \inserted{l’opération de} l’élément de $\unsure{\Pi|\pi}_1(X)$, \uline{frobénius} en la place \deleted{\unsure{f}} réelle - !]

% slanted
\margin{à vérifier}

% archive sheet number: 202
\scanpage{204}
ALGÈBRES de LIE SEMI-SIMPLES \mspage{(3)}

Tribu n° 59

et $\mathrm{q}:\mathrm{S}\longrightarrow\mathrm{W}$, et si on suppose que $\mathrm{q}(\mathrm{s})\mathrm{B}(\mathrm{v},\mathrm{v}')=\mathrm{B}(\mathrm{r}(\mathrm{s})\mathrm{v},\mathrm{v}')+\mathrm{B}(\mathrm{v},\mathrm{r}'(\mathrm{s})\mathrm{v}')$, (resp. $=\mathrm{B}(\mathrm{r}(\mathrm{s})\mathrm{v},\mathrm{r}'(\mathrm{s})\mathrm{v}')$, pour tout $\mathrm{s}$ de $\mathrm{S}$, alors $\mathrm{B}$ applique $\mathrm{V}_{\lambda}\times\mathrm{V}'_{\mu}$ dans $\mathrm{W}_{\lambda+\mu}$, (resp. dans $\mathrm{W}_{\lambda\mu}$) ; idem pour les indices supérieurs.

Ce n° absorde un certain nombre de choses du texte.

\uline{n° 1} - Vider p. 1, les l.9-7 bas, (c'est faux). La prop. 1, le lemme, la prop. 2, a) et la prop. 3 (a,c) vont au dépotoir, (noter que la prop. 3, c) absorbe le lemme. Ici le b) de la prop. 3 sera un cor. de la prop. 2, (réduite à son b)).

Dans la prop. 2, b), un erratum à Alg. VII permettra de parler des valeurs propres de $\mathrm{x}_{\mathrm M}$ même si elles ne sont pas dans $\mathrm{k}$. Pas besoin de supposer $\underline{\mathrm h}$ nilpotente à la p. 3, l.4 ; raisonner directement en utilisant Alg. VII. Dans (i), ne pas faire intervenir le quotient simple, (y'en a pas en caract. $\mathrm{p}$). Pour (ii), utiliser le dépotoir : la sous-algèbre engendrée par $\mathrm{x}$ opère sur $\mathrm{M}$, et sur $\underline{\mathrm h}$ par représ. adjointe ; on a $\underline{\mathrm h}\times\mathrm{M}\longrightarrow\mathrm{M}$.

La ``descente'' du bas de la p. 3 sera faite en Alg. VII ; donc un rappel suffit ici ; prendre une extension galoisienne qqconque au lieu de $\bar{\mathrm{k}}$.

Dans la prop. 4, à la l.13 p. 5, dire : \ldots\ est $\underline{\mathrm{g}}$-invariante, (identité de Jacobi), donc $\underline{\mathrm{h}}$-invariante. Dire ``$\underline{\mathrm{h}}$-module trivial'' à la l.15. Dire ``appliqué sur $0$ par la forme ci-dessus'' à la l.16. Plus bas la sommation est incorrecte : sommer sur la moitié des $\alpha$. Vider la remarque du bas de la p. 5.

On vide la prop. 5 qui ne sert à rien. Une remarque donnera la nilpotence de $\operatorname{Ad}(\mathrm{x})$ sous les conditions de la prop. 4.

Remonter ici le résultat suivant, (venant du \S\ 5, n° 3) : soient $\mathrm{G}$ un groupe analytique, $\mathrm{n}$ le rang de son algèbre de Lie ; alors il existe un sous-groupe ouvert $\mathrm{U}$ de $\mathrm{G}$ tel que, pour tout $\mathrm{x}\in\mathrm{U}$, $\operatorname{Ad}(\mathrm{x})$ admette $1$ comme valeur propre de

% archive sheet number: 203
\scanpage{205}
Ainsi, $\unsure{\Pi|\pi}_1(X)$ n’opère évidemt pas sur $\deleted{\unsure{\bar{\xi}}}\unsure{\xi|\zeta}_{x,\mathbb{Z}}$, ni \unsure{$=$|même} sur $\xi_{x,\mathbb{R}}$, mais l’élément $f_\infty$ de $\Pi_1(X)$ (frobénius infini) opère sur $\xi_{x,\mathbb{Z}}$ et à fortiori sur $\xi_{x,\mathbb{R}}$ ; on le prolonge \uline{linéairement} en $\tau=f_\infty$ \ldots.

\uline{N.B.} La \uline{conjecture de Hodge} interprète la filtration \uline{induite} sur la coh. ordinaire transcendante $\xi_{x,\mathbb{Z}}\otimes_{\mathbb{Z}}\mathbb{Q}\inserted{=\xi_{x,\mathbb{Q}}}$ par la filtration arithmétique ($=$ motivique, par le niveau), en termes de la position \unsure{par.} à la graduation [i.e. en termes de la filtration de Hodge --- non motivique --- et du $f_\infty$]. Elle s’exprime en disant que

\deleted{$\unsure{F|\xi}$ de niveau $\leq\unsure{p|\rho}\iff$}

\[
\begin{aligned}
&\text{La \uline{partie de niveau}}\ \leq p\quad\text{\uline{dans}}\quad\xi_{x,\mathbb{Q}}(F)\\
&\text{\uline{n’est autre que}}\quad \xi_{x,\mathbb{Q}}(F)\cap\xi_{x,\infty}(F)^{\unsure{(p)|p|(\rho)}}\cap f_\infty\xi_{x,\infty}(F)^{(p)}.
\end{aligned}
\]

En particulier, ce dernier sous-espace de $\xi_{x,\mathbb{C}}(F)$ est \uline{motivique} (provient d’un \deleted{\unsure{un}} sous-motif $F^{(p)}$ de $F$), alors que $\xi_{x,\infty}(F)^{(p)}$ et $f_\infty\xi_{x,\infty}(F)^{(p)}$ n’est pas motivique \ldots. \deleted{Si $p=0$, cela donne l’intersection des $2$ réseaux ``\unsure{entiers|rationnels|niveau}''}

% archive sheet number: 204
\scanpage{206}
Tribu n° 59

DÉCISIONS sur les ALGÈBRES de LIE SEMI-SIMPLES

(n° 400)

-:-:-:-

Le nouveau titre sera ``algèbres de Lie semi-simples déployées'' ; les autres seront dites ``repliées''. Le prochain rédacteur suivra d'assez près la présente rédaction et y ajoutera :

a) Les modules infinis $\mathrm{W}(\lambda)$, (cf. \S\ 4), avec formules de Harish-Chandra, et résultats dudit sur les opérateurs différentiels invariants ;

b) Jacobson-Mort aux Off ;

c) Les algèbres déployées classiques, et leurs représentations fondamentales, (ne parler que des représentations d'algèbres) ; montrer les autom. de trialité de $\mathrm{D}_4$ ; les algèbres exceptionnelles iront en exercices.

D'autre part Cartier rédigera un tapis arithmétique sur les $\mathrm{Z}$-structures. Les autres compléments dont parlent les commentaires du n° 400 sont, ou bien vidés, ou bien incorporés dans les \S\S\ déjà rédigés.

\uline{Annexe} I \uline{des applications polynomiales}.

Noter que l'annexe II est traitée avec la rédaction des racines.

\uline{n° 1} - Dans l'énoncé de la prop. 1 ajouter ``En particulier toutou vert non vide est dense'', (ce qui supprime le cor.) ; prendre la démonstration par références.

\uline{n° 2} - A la p. A3, l.7, ajouter ``pour les topol. de Zariski''. Juste après intervertir définition et critère pour les applic. dominantes. A la l.7 bas (p. A3), remplacer $\mathrm{W}$ par $\mathrm{A}_{\mathrm W}$.

% archive sheet number: 205
